% TOP_PROBLEM: {"id": 1101203, "title": "Questions in Geometric Group Theory — Q 12.3", "category_id": 17, "category": {"id": 17, "name": "group_theory", "display_name": "Group Theory", "description": "Problems about groups, group actions, representations, and related algebraic structures.", "slug": "group-theory", "order_index": 17, "created_at": "2026-07-31T15:26:25.671Z"}, "status": "partially_solved", "problem_number": "AMR-010-1203", "set_id": 13, "statusQualification": "Part (a) is resolved by Kent--Leininger; the definition preserves it as context and poses both remaining alternatives in part (b)."}
% FIRST_POSTED: 2026-10-01
% ENTRY_KIND: statement_only
% UnsolvedMath ID 1101203; source catalogue code AMR-010-1203
% !TeX program = lualatex
% Definition and sources only; this file is not an LLM solution attempt.
\documentclass[11pt,a4paper]{article}
\usepackage[margin=25mm]{geometry}
\usepackage{fontspec}
\setmainfont{lmroman10-regular.otf}[BoldFont=lmroman10-bold.otf,
 ItalicFont=lmroman10-italic.otf,BoldItalicFont=lmroman10-bolditalic.otf]
\usepackage{amsmath,amssymb,mathtools}
\usepackage{unicode-math}
\setmathfont{latinmodern-math.otf}
\AtBeginDocument{\let\setminus\smallsetminus}
\usepackage{xurl,hyperref}
\hypersetup{colorlinks=true,urlcolor=blue}
\setcounter{secnumdepth}{0}
\setlength{\parindent}{0pt}
\setlength{\parskip}{5pt}
\setlength{\emergencystretch}{3em}
\begin{document}
\section*{Questions in Geometric Group Theory — Q 12.3}\label{1101203}
{\small UnsolvedMath ID 1101203 · AMR-010-1203 · Group Theory · AMR Open Problem Lists}
\subsection{Definitions and mathematical statement}
Let $\Sigma_g$ denote a closed connected orientable surface of genus $g\geq2$. Its mapping class group $\operatorname{Mod}(\Sigma_g)$ consists of orientation-preserving homeomorphisms modulo isotopy. A pseudo-Anosov class has a representative preserving a pair of transverse singular measured foliations, stretching one transverse measure by $\lambda>1$ and contracting the other by $\lambda^{-1}$.

The source's part (a) asks for an injective homomorphism from some $\pi_1(\Sigma_g)$ to some $\operatorname{Mod}(\Sigma_h)$, $g,h\geq2$, whose nonidentity image elements are all pseudo-Anosov. This existence assertion is answered affirmatively by Kent--Leininger [S3], and is recorded as context rather than as an unresolved target.

Part (b) asks whether there exist $g,h\geq2$ and a smooth surface bundle
\[
 \Sigma_h\longrightarrow M\longrightarrow\Sigma_g
\]
whose fundamental group is word-hyperbolic. A bundle is locally a product with transition maps diffeomorphisms of the fiber; its total space is a closed four-manifold. Word-hyperbolicity means that some Cayley graph of the fundamental group has uniformly thin geodesic triangles.

The stronger geometric version asks whether such an $M$ can admit a Riemannian metric of constant sectional curvature $-1$. Both genera and the bundle may vary. The existence of the purely pseudo-Anosov representation in (a) does not itself assert either conclusion in (b). In particular, absence of a subgroup $\mathbb Z^2$ alone is not the definition of hyperbolicity. A metric satisfying the stronger target would imply the group-theoretic target.
\subsection{Short English statement}
Can a closed surface bundle over a closed surface have hyperbolic fundamental group, or even a real hyperbolic metric?
\subsection{Sources}
\begin{itemize}
\item[S1] ULAM.AI and the UnsolvedMath Contributors, supplied problems.json, record 1101203 (AMR-010-1203), AMR Open Problem Lists. Catalogue identity and source transcription; CC BY 4.0.
\url{https://huggingface.co/datasets/ulamai/UnsolvedMath}.
\item[S2] Bestvina - Questions in Geometric Group Theory (pdf) (2004), Question 12.3 (PDF page 19). Original source indexed by AMR-010-1203.
\url{https://www.math.utah.edu/~bestvina/eprints/questions-updated.pdf}.
\item[S3] Autumn E. Kent and Christopher J. Leininger, Atoroidal surface bundles. Purely pseudo-Anosov surface subgroups.
\url{https://annals.math.princeton.edu/articles/22768}.
\end{itemize}
\end{document}
